% page 428 of the facsimile (image p-427 of the scan)
% block 157.5 x 259.3 mm ; left margin 8.0 mm
\pgc{-12.700mm}{0.000mm}{
\l{}
\l{}
\l{}
\l{\cel{3}420}
\l{}
\l{\sou{papagayo}.\cel{2}(\sou{zool.)}\cel{3}Ucelo ek la klaso \textquotedbl{}klimeri\textquotedbl{}, di qua la}
\l{\cel{3}plumaro esas, che preske omni, bunte verda, reda, e c.,}
\l{\cel{3}kun beko harda e kurvatra, qua imitas facile la voco ho-}
\l{\cel{3}mala, la krii da la animali cetera. - L.\cel{1}\sou{psittacus}. DIRS.}
\l{}
\l{\sou{papavero}. (\sou{bot.)}\cel{3}Planto narkotigiva. - L.\cel{1}\sou{papaver rhoeas}.}
\l{\cel{3}- EFIL.}
\l{}
\l{\sou{papayo.}\cel{2}\sou{(bot.)}\cel{2}Frukto de arboro ek la regioni tropikala -}
\l{\cel{3}qua aspektas e saporas kelke quale melono. - L.\cel{1}\sou{carica}}
\l{\cel{3}\sou{papaya.}\cel{3}DEFIRS.}
\l{}
\l{\sou{papero}. - Folio por skribar od imprimar, por envelopar, e c.,}
\l{\cel{3}facita ek shifoni, papera rezidui, diversa materii ve-}
\l{\cel{3}jetantala fibroza (esparto, maizo, e c.) quin onu pastigas,}
\l{\cel{3}extensas, sekigas. - DEFS.}
\l{}
\l{\sou{papilo}.\cel{2}(\sou{anat.})\cel{2}Mikra saliajo, sub epidermo o sur la surfaco}
\l{\cel{3}di ula mukozi, ek branchifuri nerva, vaskula, qua pro-}
\l{\cel{3}duktas la sentiveso di la pelo.\cel{2}- DEFIS.}
\l{}
\l{\sou{papiliono}. (\sou{zool.)}\cel{3}Insekto qua havas quar ali, kolorizita}
\l{\cel{3}ye squami dina quale polvo. - EFL.}
\l{}
\l{\sou{papiro}. Folieto deprenata de sorto di kano qua kreskas en}
\l{\cel{3}Egiptia, India, e c., e quan onu preparis (en la epoki}
\l{\cel{3}antiqua) por la skribo. -\cel{2}DEFIRS.}
\l{}
\l{\sou{paplo}. I. Farino o fekulo quan onu lasas boliar en lakto,}
\l{\cel{3}aquo, por facar de lu nutrivo destinata precipue a la}
\l{\cel{3}infanti. - II. (\sou{metaf}.)Shifoni boliita e reduktita a pasto,}
\l{\cel{3}de qua facesos papero, kartono, e c. - EIS.}
\l{}
\l{\sou{papriko.(bot.})\cel{2}Varietato di pimento, tre uzata en Hungaria.}
\l{\cel{3}- L.\cel{1}\sou{capsicum}.\cel{3}DFL.}
\l{}
\l{\sou{papulo}.\cel{2}I. (\sou{patol.}) Lezuro di la pelo, karakterizata da}
\l{\cel{3}saliajo (rozea, reda o bruna) di la pelo. - II. Protube-}
\l{\cel{3}ranco mola e plena ye liquido aquatra, sur la epidermo di}
\l{\cel{3}ula planti. - EFIS.}
\l{}
\l{\sou{par-}\cel{1}.\cel{2}Prefixo qua signifikas : \textquotedbl{}komplete\textquotedbl{}, kun la radiki di}
\l{\cel{3}verbi.}
\l{}
\l{\sou{paro}. I. Duo de kozi od animali, sama-speca, destinita irar}
\l{\cel{3}kune. - II. La maskulo e la femino di animali sama-speca,}
\l{\cel{3}un de singlu. -}
\l{}
\l{\sou{para-}\cel{4}Prefixo qua signifikas \textquotedbl{}qua protektas kontre\textquotedbl{},}
\l{\cel{3}\textquotedbl{}shirmilo kontre\textquotedbl{}.}
\l{}
\l{\sou{parabolo}. I. (\sou{retoriko})\cel{2}Alegorio sub qua celesas docivo eti-}
\l{\cel{3}kala, morala. - II. (\sou{geom.)}\cel{3}Kurvo plana di la grado du-}
\l{\cel{3}esma, qua rezultas de la seciono di kono dsa plano paralela}
\l{\cel{3}ye la latero. - DEFIRS.}
}
